% 3-8-ensemble.tex
%  Budget: 3pp.  Static ensemble first, then regime-aware.
%  Source map: forecast combination as a standard accuracy gain in EPF
%   -> jiang_review_2018 ; the published ensembles this thesis is measured
%   against -> lago_forecasting_2021 .
%  Constants verified against configs/evaluation.yaml and data/raw/DE.csv:
%   threshold 62.6989 = train mean 37.6876 + 1.5 * std 16.6742, computed on
%   prices up to and including 2015-01-04; validation split 37/357;
%   test split 77/728. Oracle bound from reports/tables/oracle_bound.csv:
%   3.5580 (frozen seed-42 LSTM member set).
%  MANDATORY DISCLOSURES (HANDOFF 16A/A6 + A8, and section 5.3):
%   (a) k=1.5 chosen by a VALIDATION-ONLY rule -- test never consulted;
%   (b) the regime label uses the PREVIOUS day's realized prices;
%   (c) ensemble weights were fitted on the members' TUNING window.
%  FORBIDDEN: 84.04 as a threshold; the label جهش / "spike" for the regime.
%  The fitted weight vectors are NOT in any frozen table -> flagged, not invented.

\section{مدل ترکیبی}
\label{sec:3-8}

\subsection*{چرا ترکیب}

هر یک از چهار مدل برازشیِ این پژوهش خطای متفاوتی مرتکب می‌شود. مدل خطی در شرایطی
می‌لغزد که مدل درختی از پس آن برمی‌آید و بالعکس. ترکیب وزن‌دار، همین ناهم‌بستگیِ
خطاها را به دقت تبدیل می‌کند و در ادبیات پیش‌بینی قیمت برق از راه‌های شناخته‌شده‌ی
بهبود دقت است \cite{jiang_review_2018}؛ گونه‌های ترکیبیِ خودِ مقاله‌ی مرجع نیز بهترین
اعداد آن مقاله‌اند \cite{lago_forecasting_2021}.

دو مدل ترکیبی در اینجا معرفی می‌شود: نخست ترکیب ایستا، سپس ترکیب رژیم‌محور که سهم
مشارکت پژوهش حاضر است.

\subsection*{ترکیب ایستا}

ترکیب ایستا یک بردار وزنِ ثابت بر چهار مدل برازشی \lr{—} \lr{SARIMAX}،
\lr{LEAR-LASSO}، \lr{LightGBM} و \lr{LSTM} \lr{—} اعمال می‌کند. مدل ساده در ترکیب
شرکت داده نمی‌شود، چون نقش آن مرجع است، نه عضو.

وزن‌ها با کمینه‌سازیِ خطای مطلق میانگین بر پنجره‌ی برازشِ وزن به دست می‌آیند، با دو قید
که ترکیب را محدب نگه می‌دارند:

\begin{equation}
\hat{w} = \arg\min_{w}
\frac{1}{N}\sum_{i=1}^{N}\left| p_i - \sum_{m=1}^{M} w_m \hat{p}_{m,i} \right|
\quad \text{به شرط} \quad
w_m \geq 0, \;\; \sum_{m=1}^{M} w_m = 1
\end{equation}

که در آن $M$ شمار مدل‌های عضو، $\hat{p}_{m,i}$ پیش‌بینیِ مدل $m$ در نقطه‌ی $i$، و
$w_m$ وزن آن مدل است. قید نامنفی‌بودن و قید جمع‌شدن به یک، هر دو عامدانه‌اند: بدون
آن‌ها بهینه‌ساز می‌تواند وزن‌های بزرگ و مخالف‌علامت پیدا کند که بر پنجره‌ی برازش عالی
عمل می‌کنند و بیرون از آن ناپایدارند.

پیاده‌سازی در برابر شکست خاموش محافظت شده است: اگر بهینه‌ساز همگرا نشود یا وزن‌هایی
بازگرداند که به یک جمع نمی‌شوند، اجرا با خطا متوقف می‌شود، به‌جای آنکه وزن‌های
یکنواخت را در سکوت جایگزین کند و جدول‌های پس از آن همچنان برچسب «وزن‌های بهینه» را
حمل کنند.

\textbf{افشای نخست.} وزن‌ها بر پنجره‌ی \emph{تنظیمِ} خودِ اعضا برازش شده‌اند، یعنی
بازه‌ای که آن اعضا ابرپارامترهایشان را از آن گرفته‌اند. این را باید گفت. آنچه آن را
مهار می‌کند این است که مقایسه‌ای که این وزن‌ها تغذیه می‌کنند \lr{—} یعنی رژیم‌محور در
برابر ایستا \lr{—} همین مزیت را به‌طور یکسان به هر دو طرف می‌دهد، پس تفاوتِ میان آن
دو از این بابت مخدوش نمی‌شود. آنچه مخدوش می‌شود مقایسه‌ی ترکیب با تک‌مدل‌هاست، که با
احتیاط بیشتری باید خوانده شود.

\subsection*{تعریف رژیم بازار}

\cref{sec:3-3-3} نشان داد که نوسان در این سری خوشه‌ای است: دوره‌های آرام و پرتلاطم کنار هم
می‌نشینند. اگر رفتار بازار در این دو حالت به‌قدر کافی متفاوت باشد، بردار وزنی که در
حالت آرام بهینه است لزوماً در حالت پرتنش بهینه نیست. ترکیب رژیم‌محور بر همین فرض
استوار است.

دو رژیم تعریف می‌شود \lr{—} \textbf{آرام} و \textbf{پرتنش} \lr{—} و مرز آن‌ها یک
آستانه‌ی قیمتی است:

\begin{equation}
\tau = \mu_{\mathrm{train}} + k \cdot \sigma_{\mathrm{train}}
\end{equation}

که در آن $\mu_{\mathrm{train}}$ و $\sigma_{\mathrm{train}}$ میانگین و انحراف معیارِ
قیمت بر بخشی از داده‌ی آموزش، و $k$ ضریبی است که باید انتخاب شود. مقادیر عددی چنین
است: میانگین \lr{37.6876} و انحراف معیار \lr{16.6742} بر قیمت‌های تا پایان ۴ ژانویه‌ی
۲۰۱۵، که با $k$ برابر \lr{1.5} آستانه‌ی \lr{62.6989} یورو بر مگاوات‌ساعت را می‌دهد.

سه ویژگیِ این تعریف را باید صریح گفت، چون هر سه پرسش‌هایی را پیشاپیش می‌بندند که
حتماً پرسیده می‌شوند.

\textbf{یکم: مبنای محاسبه‌ی آستانه پیش از هر پنجره‌ی آلوده‌کننده تمام می‌شود.} برش ۴
ژانویه‌ی ۲۰۱۵ اکیداً پیش از هر دو پنجره‌ای قرار دارد که می‌توانستند آن را آلوده کنند:
پنجره‌ی تنظیم ابرپارامتر و پنجره‌ی برازش وزن، که هر دو از ۲۰۱۵ آغاز می‌شوند. آستانه
بنابراین نه از دوره‌ی آزمون و نه از پنجره‌های تنظیم، اطلاعی ندارد. این مقدار در
آزمون‌های پروژه در برابر داده‌ی خام سنجیده می‌شود، یعنی خاصیت وارسی می‌شود و نه صرفاً
در یک توضیح ادعا.

\textbf{دوم: $k$ با قاعده‌ای صرفاً بر پایه‌ی داده‌ی اعتبارسنجی انتخاب شده است.} قاعده
این بود: بزرگ‌ترین $k$ از میان مقادیر \lr{3.0}، \lr{2.5}، \lr{2.0} و \lr{1.5} که به
ازای آن هر دو رژیم دست‌کم بیست روز در پنجره‌ی اعتبارسنجی داشته باشند. مقدار \lr{2.0}
تنها ده روز پرتنش باقی می‌گذاشت و \lr{1.5} سی‌وهفت روز؛ پس \lr{1.5} انتخاب شد.
\textbf{دوره‌ی آزمون هرگز در این انتخاب مشورت نشد.} انتخاب آستانه بر پایه‌ی رفتار آن
روی دوره‌ی آزمون، گزینشِ درون‌نمونه‌ای می‌بود و کل نتیجه‌ی \cref{sec:4-5-1} را بی‌اعتبار
می‌کرد.

تقسیم حاصل چنین است: در پنجره‌ی اعتبارسنجی ۳۷ روز پرتنش در برابر ۳۲۰ روز آرام (از مجموع ۳۵۷ مبدأ)، و در
دوره‌ی آزمون ۷۷ روز پرتنش از مجموع ۷۲۸ روز.

\textbf{سوم: برچسب رژیم از روز پیش خوانده می‌شود.} یک روز زمانی پرتنش برچسب می‌خورد
که قیمت‌های محقق‌شده‌ی \emph{روز پیش از آن} از آستانه عبور کرده باشند. این قید در کد
اعمال می‌شود، نه در توضیح. اگر برچسب از قیمت‌های خودِ روز هدف خوانده می‌شد، خودِ
برچسب حامل اطلاعاتی از آینده می‌بود و کل سازوکار به نشت اطلاعات آلوده می‌گشت.

نکته‌ی اصطلاح‌شناختی: در سطح تقریبیِ یک‌ونیم انحراف معیار، آنچه علامت‌گذاری می‌شود یک
روز \emph{بالاتر از معمول} است و نه یک جهش قیمتی. برچسب «پرتنش» عامدانه انتخاب شده
تا این تمایز حفظ شود.

\subsection*{ترکیب رژیم‌محور}

ترکیب رژیم‌محور به‌جای یک بردار وزن، دو بردار دارد: یکی برازش‌شده بر روزهای آرامِ
پنجره‌ی برازش و دیگری بر روزهای پرتنشِ همان پنجره. در زمان پیش‌بینی، برچسب رژیمِ روز
هدف \lr{—} که، مطابق بالا، از روز پیش خوانده می‌شود \lr{—} تعیین می‌کند کدام بردار
وزن به کار رود. به‌جز همین سوئیچ، هیچ تفاوت دیگری با ترکیب ایستا وجود ندارد: همان
اعضا، همان پنجره‌ی برازش، همان بهینه‌ساز و همان قیدهای محدب.

حداقلی‌بودنِ این طراحی عامدانه است. اگر ترکیب رژیم‌محور هم‌زمان اعضای دیگری می‌گرفت یا
پنجره‌ی متفاوتی به کار می‌برد، تفاوت عملکردِ آن با ترکیب ایستا دیگر به شرطی‌کردنِ
وزن‌ها قابل نسبت‌دادن نبود.

یادآوری تفکیکِ \cref{sec:3-5}: پیش‌بینی‌هایی که این وزن‌ها بر آن‌ها اعمال می‌شوند، خروجی
مدل‌های چارچوب پیش‌رونده‌اند و نه مدل تثبیت‌شده بر کل داده؛ یعنی هر پیش‌بینیِ عضو، از
مدلی می‌آید که اکیداً پیش از مبدأ خود برازش شده است. وزن‌ها نیز، همان‌گونه که گفته شد،
بر پنجره‌ای پیش از دوره‌ی آزمون برازش می‌شوند. پس هیچ مرحله‌ای از این زنجیره دوره‌ی
آزمون را نمی‌بیند.

\subsection*{وزن‌های برازش‌شده}

\cref{tab:ensemble-weights} سه بردار وزن را در کنار هم می‌آورد: بردار ایستا، و
دو بردار رژیم‌محورِ آرام و پرتنش. هر سه با همان رویه‌ی بالا و بر همان پنجره‌ی برازش به
دست آمده‌اند.

\begin{table}[htbp]
\centering
\caption{وزن‌های برازش‌شده‌ی مدل ترکیبی بر پنجره‌ی برازشِ وزن (۳۵۷ مبدأ، از ۱۲
ژانویه‌ی ۲۰۱۵ تا ۳ ژانویه‌ی ۲۰۱۶). ستون نخست وزن‌های ترکیب ایستا و دو ستون بعد
وزن‌های ترکیب رژیم‌محور است. هر ستون به یک جمع می‌شود. تقسیم رژیمی این پنجره ۳۲۰ روز
آرام و ۳۷ روز پرتنش است.}
\label{tab:ensemble-weights}
\small
\setlength{\tabcolsep}{6pt}
\begin{tabular}{lrrr}
\toprule
& & \multicolumn{2}{c}{رژیم‌محور} \\
\cmidrule(lr){3-4}
عضو & ایستا & آرام & پرتنش \\
\midrule
\lr{SARIMAX} & \lr{0.134} & \lr{0.140} & \lr{0.063} \\
\lr{LEAR-LASSO} & \lr{0.245} & \lr{0.237} & \lr{\textbf{0.375}} \\
\lr{LightGBM} & \lr{\textbf{0.382}} & \lr{\textbf{0.388}} & \lr{0.281} \\
\lr{LSTM} & \lr{0.239} & \lr{0.235} & \lr{0.281} \\
\bottomrule
\end{tabular}
\end{table}

مقایسه‌ی دو ستون آخر، خودِ سازوکار را آشکار می‌کند و از هر توضیح کلامی گویاتر است.

در رژیم آرام، بردار وزن عملاً همان بردار ایستاست؛ اختلاف‌ها در حد چند هزارم‌اند. این
انتظارپذیر است، چون روزهای آرام اکثریتِ پنجره‌ی برازش را می‌سازند و بردار ایستا نیز
عمدتاً از همان‌ها شکل می‌گیرد.

در رژیم پرتنش اما آرایش وزن‌ها به‌روشنی جابه‌جا می‌شود. سهم \lr{LEAR-LASSO} از
\lr{0.237} به \lr{0.375} می‌رسد و به بزرگ‌ترین وزن بدل می‌گردد، سهم \lr{LightGBM} از
\lr{0.388} به \lr{0.281} افت می‌کند، و سهم \lr{SARIMAX} بیش از نصف کاهش می‌یابد و از
\lr{0.140} به \lr{0.063} می‌رسد.

خوانش این جابه‌جایی چنین است: در روزهای پرتنش، بهینه‌ساز وزن را از مدل درختی به سمت
مدل خطی می‌برد. مشاهده‌ی \lr{RMSE} در \cref{sec:4-2} با همین هم‌خوان است، جایی که مدل درختی
خطاهای حدیِ بزرگ‌تری از مدل خطی نشان می‌دهد؛ درختِ تصمیم در دنباله‌های توزیع اشباع
می‌شود، حال آنکه مدل خطی در همان ناحیه برون‌یابی می‌کند. روزهای پرتنش دقیقاً همان
روزهایی‌اند که قیمت به دنباله‌ی توزیع نزدیک می‌شود.

باید تصریح کرد که این تفسیر، \emph{پس از} دیدن وزن‌ها نوشته شده است، نه پیش از آن؛
بنابراین توضیحی است بر یک الگوی مشاهده‌شده و نه فرضیه‌ای که آزموده شده باشد. آنچه
آزموده شده، تنها این است که وزن‌دهیِ رژیم‌محور دقت را بهبود می‌دهد یا نه، و پاسخ آن در
\cref{sec:4-5-1} می‌آید.

\subsection*{یک مرجع مفید: کران اوراکل}

برای آنکه بتوان قضاوت کرد بهبودِ رژیم‌محوری چه اندازه است، مرجعی لازم است: بهترین
دقتی که \emph{هر} بردار وزن ایستا می‌توانست بر دوره‌ی آزمون به دست آورد، اگر آن را با
دانستن پاسخ انتخاب می‌کردیم. این کران اوراکل بر مجموعه‌ی اعضای تثبیت‌شده برابر
\lr{3.5580} است.

مقایسه گویاست و در \cref{sec:4-5-1} بسط می‌یابد: ترکیب ایستای واقعی به \lr{3.5742}
می‌رسد، یعنی اندکی بالاتر از کران، که طبیعی است چون وزن‌هایش را بدون دیدن دوره‌ی
آزمون برازش کرده است. ترکیب رژیم‌محور به \lr{3.5569} می‌رسد \lr{—} یعنی به عددی
اندکی \emph{بهتر} از بهترین وزنِ ایستای ممکن. این خودِ ادعای \cref{sec:3-8} را می‌سازد:
بهبود از شرطی‌کردنِ وزن‌ها بر رژیم می‌آید و نه از یافتن وزن ایستای بهتر، چون هیچ وزن
ایستایی نمی‌توانست به آن برسد.

\cref{sec:4-5-1} نشان خواهد داد که این بهبود کجا رخ می‌دهد، و پاسخ \lr{—} تماماً بر روزهای
پرتنش \lr{—} همان چیزی است که این طراحی برای آن ساخته شده بود.
